\documentclass[aspectratio=169,11pt]{beamer}

\usepackage{iftex}
\ifXeTeX
  \usepackage{fontspec}
  \usepackage{xeCJK}
  \setmainfont{Times New Roman}
  \IfFontExistsTF{Microsoft JhengHei}{
    \setCJKmainfont{Microsoft JhengHei}
    \setCJKsansfont{Microsoft JhengHei}
    \setCJKmonofont{Microsoft JhengHei}
  }{
    \setCJKmainfont{SimSun}
    \setCJKsansfont{SimSun}
    \setCJKmonofont{SimSun}
  }
  \IfFontExistsTF{Consolas}{\setmonofont{Consolas}}{}
\fi
\usepackage{amsmath}
\usepackage{booktabs}
\usepackage{array}
\usepackage{graphicx}
\usepackage{siunitx}
\usepackage{xurl}
\usepackage{xcolor}

\usetheme{Madrid}
\usecolortheme{default}
\setbeamertemplate{navigation symbols}{}
\setbeamertemplate{footline}[frame number]
\graphicspath{{images/neural_network_fit_results/}}

\newcommand{\goodtext}[1]{\textcolor{green!45!black}{#1}}
\newcommand{\warntext}[1]{\textcolor{red!65!black}{#1}}
\newcommand{\notetext}[1]{\textcolor{blue!65!black}{#1}}
\newcommand{\pathtext}[1]{{\scriptsize\url{#1}}}
\newcommand{\src}[1]{\vfill{\tiny\textcolor{gray!75!black}{Source: \url{#1}}}}
\newcolumntype{L}[1]{>{\raggedright\arraybackslash}p{#1}}
\newcolumntype{R}[1]{>{\raggedleft\arraybackslash}p{#1}}

\title[NN Fitting Results]{神经网络拟合结果}
\subtitle{最新 \texorpdfstring{\(\Delta P^{0.5}\)}{Delta P\textasciicircum 0.5} Stage-3 模型、五 seed 稳定性与轨迹级诊断}
\author{Linan Jiang}
\institute{Aalto University}
\date{10 May 2026}

\begin{document}

\begin{frame}
  \titlepage
\end{frame}

\begin{frame}{为什么这版替换 4 月结果}
\small
\begin{block}{关键变化}
上游 CDF/clean 数据在 5 月流程中变得更干净且规模更大；同时特征工程从
\(\Delta P^{0.25}\) 更新到数据支持的 \(\Delta P^{0.5}\) scaling。因此，
4 月 18 日的 tuned run 不能再作为当前最佳模型，也不能与新重训模型做直接
百分比改善比较。
\end{block}

\begin{block}{本 deck 的严格口径}
\begin{itemize}
  \item headline 结果使用同一 external clean diagnostic：
  \num{795561} points，\num{33845} trajectories；
  \item 主结论使用 \(\Delta P^{0.5}\) Stage-3 pipeline 的五 seed aggregate；
  \item 图形诊断使用单个代表性最佳 seed 99，只用于展示误差结构，不替代五 seed 结论；
  \item 旧 2026-04-18 figures 仅为 archive，本版不再使用其 baseline-comparison 叙事。
\end{itemize}
\end{block}
\end{frame}

\begin{frame}{当前评估对象}
\small
\begin{tabular}{L{0.30\linewidth}L{0.60\linewidth}}
\toprule
项目 & 设置 \\
\midrule
模型族 & Stage-3 raw-CDF refinement，\(\Delta P^{0.5}\) engineered feature scaling \\
Headline protocol & external clean diagnostic，\SIrange{0}{5}{ms} \\
Headline statistic & 五 seed mean \(\pm\) seed-to-seed std / bootstrap CI \\
代表性图形 seed & seed 99，\texttt{anchor\_off}，单 seed 最低 RMSE \\
输入特征 & \(t\), tilt angle, plume count, injection duration, control backpressure \\
被移除的稀疏轴 & direct \(d_n\), \(P_{\mathrm{inj}}\), \(P_{\mathrm{ch}}\), \(\Delta P\) inputs \\
Target scaling & \(S/A,\quad A=\Delta P^{0.5}\rho_a^{-0.25}d_n^{0.5}\) \\
\bottomrule
\end{tabular}

\src{Thesis/generated/baseline_comparison_20260509/per_seed_new_mlp_eval.csv}
\end{frame}

\begin{frame}{数据规模与当前最佳结果}
\small
\begin{columns}[T,onlytextwidth]
\begin{column}{0.49\linewidth}
\begin{tabular}{L{0.55\linewidth}R{0.32\linewidth}}
\toprule
Metric & Value \\
\midrule
CSV files & 4,391 \\
Trajectories & 33,845 \\
Prediction points & 795,561 \\
Truth range & \SIrange{0.00}{95.88}{mm} \\
\midrule
5-seed RMSE & \goodtext{\SI{6.211}{mm} \(\pm\) \SI{0.283}{mm}} \\
5-seed MAE & \SI{4.587}{mm} \(\pm\) \SI{0.281}{mm} \\
5-seed Bias & \SI{-0.838}{mm} \(\pm\) \SI{0.702}{mm} \\
P95 absolute error & \SI{13.029}{mm} \(\pm\) \SI{0.411}{mm} \\
Range-normalized RMSE & 6.48\% \\
\bottomrule
\end{tabular}
\end{column}

\begin{column}{0.47\linewidth}
\begin{block}{代表性最佳单模型：seed 99}
\footnotesize
\begin{itemize}
  \item RMSE = \SI{5.948}{mm}
  \item MAE = \SI{4.383}{mm}
  \item Bias = \SI{-0.074}{mm}
  \item P95 absolute error = \SI{12.533}{mm}
  \item 1\(\sigma\)/2\(\sigma\) coverage = 0.792 / 0.985
\end{itemize}
\end{block}

\begin{block}{严格读法}
\footnotesize
seed 99 展示了当前 pipeline 能达到的最佳单次拟合状态；论文主结果仍应引用
五 seed aggregate，因为它反映训练随机性下的稳定表现。
\end{block}
\end{column}
\end{columns}

\src{Thesis/generated/baseline_comparison_20260509/new_mlp_5seed_aggregate.csv}
\end{frame}

\begin{frame}{指标含义}
\small
\begin{tabular}{L{0.23\linewidth}L{0.67\linewidth}}
\toprule
指标 & 含义 \\
\midrule
RMSE & \(\sqrt{\frac{1}{n}\sum_i(\hat y_i-y_i)^2}\)，对大误差敏感，是主要 accuracy 指标。 \\
MAE & \(\frac{1}{n}\sum_i|\hat y_i-y_i|\)，表示典型毫米误差。 \\
Bias & \(\frac{1}{n}\sum_i(\hat y_i-y_i)\)，保留符号，检验系统性高估/低估。 \\
P95 abs. err. & \(|\hat y-y|\) 的 95th percentile，报告尾部风险。 \\
Coverage & \(\Pr(|y-\mu|\leq k\sigma)\)。校准 Gaussian 的 1\(\sigma\)/2\(\sigma\) nominal 为 0.683/0.954。 \\
\bottomrule
\end{tabular}

\vspace{0.6em}
\footnotesize
因此，本 deck 不再说“比 4/16 baseline 改善多少”，而是报告当前 clean protocol
下的绝对误差、seed 稳定性、尾部误差和不确定度覆盖。
\end{frame}

\begin{frame}{Predicted vs. Measured Penetration}
\begin{columns}[T,onlytextwidth]
\begin{column}{0.57\linewidth}
  \centering
  \includegraphics[width=\linewidth,height=0.73\textheight,keepaspectratio]{latest_pred_vs_actual_seed99.png}
\end{column}
\begin{column}{0.39\linewidth}
\footnotesize
\begin{block}{读图方式}
\begin{itemize}
  \item 每个点是 external clean set 中一个 time sample。
  \item \(y=x\) 线表示预测穿深等于测量穿深。
  \item seed 99 的 bias 接近零，说明全局系统性低估在最新模型中已显著减弱。
  \item 高穿深区域仍存在尾部散点，需要与 P95 和 trajectory review 一起读。
\end{itemize}
\end{block}
\end{column}
\end{columns}
\src{MLP/eval/rmse_eval_clean_20260509_215303_newdp_seed99_points/points.csv}
\end{frame}

\begin{frame}{Residual Distribution 与 Calibration}
\begin{columns}[T,onlytextwidth]
\begin{column}{0.56\linewidth}
  \centering
  \includegraphics[width=\linewidth,height=0.72\textheight,keepaspectratio]{latest_residual_histogram_seed99.png}
\end{column}
\begin{column}{0.40\linewidth}
\footnotesize
\begin{block}{Residual 定义}
\[
e = \hat y_{\mathrm{pen}} - y_{\mathrm{pen}}.
\]
\end{block}

\begin{itemize}
  \item seed 99 bias = \SI{-0.074}{mm}，接近无全局偏置。
  \item 1\(\sigma\) coverage = 0.792，高于 Gaussian nominal 0.683。
  \item 2\(\sigma\) coverage = 0.985，高于 Gaussian nominal 0.954。
  \item 该 clean protocol 上模型偏保守，但作为 screening surrogate 是可接受方向。
\end{itemize}
\end{column}
\end{columns}
\end{frame}

\begin{frame}{五 seed 稳定性}
\begin{columns}[T,onlytextwidth]
\begin{column}{0.58\linewidth}
  \centering
  \includegraphics[width=\linewidth,height=0.66\textheight,keepaspectratio]{latest_per_seed_rmse.png}
\end{column}
\begin{column}{0.38\linewidth}
\scriptsize
\begin{block}{核心读法}
\begin{itemize}
  \item 5 个 seed 的 RMSE spread 约 \(\pm\SI{0.283}{mm}\)。
  \item seed 99 是最佳单 seed；seed 2024 选择了不同 ablation \texttt{raw\_reliable\_no\_kd}。
  \item 若 thesis 需要严格 same-ablation 稳定性，应单独报告四个 \texttt{anchor\_off} seeds。
\end{itemize}
\end{block}

\begin{block}{本 deck 结论}
\footnotesize
用五 seed mean 作为主结论，用 seed 99 图形解释误差结构；两者不能混成同一个统计对象。
\end{block}
\end{column}
\end{columns}
\src{Thesis/generated/baseline_comparison_20260509/per_seed_new_mlp_eval.csv}
\end{frame}

\begin{frame}{Per-Folder Error 与 Coverage}
\begin{columns}[T,onlytextwidth]
\begin{column}{0.50\linewidth}
  \centering
  \includegraphics[width=\linewidth,height=0.72\textheight,keepaspectratio]{latest_per_folder_rmse_seed99.png}
\end{column}
\begin{column}{0.48\linewidth}
  \centering
  \includegraphics[width=\linewidth,height=0.72\textheight,keepaspectratio]{latest_per_folder_coverage_seed99.png}
\end{column}
\end{columns}

\vspace{0.1em}
\footnotesize
Nozzle7/5/0 等 folder 误差较低；BC20241014\_HZ\_Nozzle3 仍是最困难组，
同时 1\(\sigma\) coverage 偏低，是当前模型的主要 trajectory-level 风险点。

\src{MLP/eval/rmse_eval_clean_20260509_215303_newdp_seed99_points/per_folder.csv}
\end{frame}

\begin{frame}{代表性轨迹：低误差与高误差}
\begin{columns}[T,onlytextwidth]
\begin{column}{0.49\linewidth}
  \centering
  \includegraphics[width=\linewidth,height=0.67\textheight,keepaspectratio]{latest_best_traj_seed99.png}

  \footnotesize
  Nozzle0 T3 plume 2：RMSE \SI{0.36}{mm}，1\(\sigma\) coverage 1.00。
\end{column}
\begin{column}{0.49\linewidth}
  \centering
  \includegraphics[width=\linewidth,height=0.67\textheight,keepaspectratio]{latest_worst_traj_seed99.png}

  \footnotesize
  BC20241014\_HZ\_Nozzle3 T6 plume 4：RMSE \SI{38.77}{mm}，1\(\sigma\) coverage 0.11。
\end{column}
\end{columns}
\end{frame}

\begin{frame}{结论}
\small
\begin{enumerate}
  \item \goodtext{当前最佳可报告结果}：\(\Delta P^{0.5}\) Stage-3 MLP 五 seed external-clean RMSE 为 \SI{6.211}{mm} \(\pm\) \SI{0.283}{mm}。
  \item \goodtext{最新单 seed 诊断}：seed 99 达到 RMSE \SI{5.948}{mm}、MAE \SI{4.383}{mm}、bias \SI{-0.074}{mm}，说明新 pipeline 的全局低估问题已显著减弱。
  \item \notetext{不确定度读法}：clean set 上 1\(\sigma\)/2\(\sigma\) coverage 为 0.774/0.981（五 seed mean），偏保守而非 under-coverage。
  \item \warntext{剩余风险}：Nozzle3 仍有明显 OOD/geometry-dependent failure mode；trajectory-level review 仍然必要。
\end{enumerate}

\begin{block}{论文口径}
不要再把 5 月重训模型与 4 月旧模型做直接改进百分比比较；它们来自不同 clean 数据规模和不同 feature-scaling protocol。
\end{block}
\end{frame}

\begin{frame}{文件索引}
\small
\begin{tabular}{L{0.24\linewidth}L{0.68\linewidth}}
\toprule
内容 & 文件 \\
\midrule
Slides source & \pathtext{Thesis/latex/neural_network_fit_results_slides_zh.tex} \\
Compiled PDF & \pathtext{Thesis/pdf/neural_network_fit_results_slides_zh.pdf} \\
Latest figures & \pathtext{Thesis/images/neural_network_fit_results/latest_*.png} \\
5-seed aggregate & \pathtext{Thesis/generated/baseline_comparison_20260509/new_mlp_5seed_aggregate.csv} \\
Representative seed eval & \pathtext{MLP/eval/rmse_eval_clean_20260509_215303_newdp_seed99_points/} \\
Asset manifest & \pathtext{Thesis/images/neural_network_fit_results/latest_assets_manifest.csv} \\
\bottomrule
\end{tabular}
\end{frame}

\end{document}
